\section*{Methods}

\subsection*{Task and latent learning state}
\label{m:task}

Each trial presents a stimulus $\bm x_t\in[0,1]^4$, followed by a category response $y_t$ and feedback $r_t$. The four coordinates encode normalized feature lengths. Condition~1 has two categories and binary feedback; condition~2 has four categories and binary feedback. Condition~3 has four categories grouped into two families. Its feedback distinguishes the correct species ($r_t=1$), the correct family but incorrect species ($r_t=0.5$), and the incorrect family ($r_t=0$). Participants in this condition know the family--species structure and learn the response-label pairings.

\label{m:state}
The model represents a classification rule as a partition of stimulus space into labeled category regions. A finite catalogue $\mathcal H$ contains rules based on feature thresholds, comparisons and conditional judgments. The learner considers an active subset $\mathcal A_t$, with size bounded by the participant's capacity $M^{\max}$. Each active rule $h$ has a normalized belief $\pi_t^-(h)$ and a prediction precision $\beta_t(h)$. The learner starts with one uniformly sampled rule, belief~1 and initial precision $\beta_0$. Superscripts $-$ and $+$ denote beliefs before choice and after feedback. Participant indices are omitted; parameters remain fixed within a participant while these states evolve.

Perception and choice precede feedback on each trial. Feedback then updates beliefs and precisions, after which hypothesis updating prepares the candidate set for the next trial. Thus, changes induced by current feedback influence subsequent choices.

\draftslot{Participant recruitment, task administration, exclusions and ethics information to be added.}

\subsection*{Perception (P): representing stimulus uncertainty}
\label{m:perception}

Perception (P) introduces feature-specific uncertainty calibrated from independent perceptual measurements. All active rules receive the same perceived stimulus $\widetilde{\bm x}_t$ within a learning trajectory. A rule predicts categories according to the distance $d_{t,h,c}$ from this stimulus to its region for category $c$:
\begin{equation}
p_{t,h}(c)=
\frac{\exp[-\beta_t(h)d_{t,h,c}]}
{\sum_{a=1}^{C}\exp[-\beta_t(h)d_{t,h,a}]},
\label{eq:category-prediction}
\end{equation}
where $C$ is the number of categories. Distance is the shortest Euclidean distance to the category region, and equals zero inside it. Greater precision makes category probabilities more sensitive to these distance differences.

\subsection*{Choice readout (R): mapping beliefs to choices}
\label{m:choice}
Choice readout (R) converts rule beliefs into contributions to the response. It also represents persistence in rule use: the previously dominant rule receives an advantage controlled by inertia $\lambda$. The choice weights and resulting category probabilities are
\begin{align}
w_t(h)&=\frac{\pi_t^-(h)\exp[\lambda\ind(h=h_{t-1}^{\star})]}
{\sum_{v\in\mathcal A_t}\pi_t^-(v)\exp[\lambda\ind(v=h_{t-1}^{\star})]},
\label{eq:choice-weights}\\
P_t(c\mid z_t)&=\sum_{h\in\mathcal A_t}w_t(h)p_{t,h}(c),
\label{eq:choice-mixture}
\end{align}
where $\ind$ is the indicator function and $z_t$ denotes the internal state, including the perceived stimulus. The current dominant rule $h_t^{\star}$ has the largest choice weight and supplies inertia on the next trial. When there is no preceding active dominant rule, weights equal beliefs. Distributed weights allow several rules to influence choice; concentrated weights allow one rule to dominate. Inertia can therefore preserve an established policy while evidence for alternatives accumulates.

\subsection*{Evidence memory (M): updating rule beliefs}
\label{m:memory}

Evidence memory (M) evaluates how well every active rule predicts the feedback received after the chosen response. For binary feedback, the likelihood $\ell_t(h)$ and updated belief are
\begin{align}
\ell_t(h)&=r_t p_{t,h}(y_t)+(1-r_t)[1-p_{t,h}(y_t)],
\label{eq:feedback-likelihood}\\
\pi_t^+(h)&=
\frac{[\pi_t^-(h)]^{\gamma}\ell_t(h)}
{\sum_{v\in\mathcal A_t}[\pi_t^-(v)]^{\gamma}\ell_t(v)}.
\label{eq:evidence-memory}
\end{align}
Correct feedback supports probability assigned to the chosen category; incorrect feedback supports the remaining categories. The retention parameter $0\leq\gamma\leq1$ controls how much previous evidence persists. Values near~1 preserve accumulated belief differences, while smaller values reduce these differences before new evidence enters. Normalization produces relative support within the current candidate set. All active candidates undergo this update, including those contributing little to choice.

\subsection*{Precision updating (B): adjusting prediction precision}
\label{m:precision}
Precision updating (B) calibrates the concentration of each rule's predictions using the feedback likelihood. It changes precision in the direction that increases the likelihood of the observed feedback under that rule. To express this update across feedback conditions, let $a_{t,h}(c)\in[0,1]$ denote feedback compatibility. It is the probability of the observed feedback given category $c$, the chosen response and the rule's existing knowledge. For binary feedback,
\begin{equation}
a_{t,h}(c)=
\begin{cases}
\ind(c=y_t),&r_t=1,\\
\ind(c\neq y_t),&r_t=0.
\end{cases}
\label{eq:feedback-compatibility}
\end{equation}
Correct feedback selects the chosen category, whereas incorrect feedback retains all other categories as possible explanations. Condition~3 averages these compatibilities over the learner's pairing beliefs, as defined below.

Write $p_{t,h}(c;\beta)$ for Eq.~\eqref{eq:category-prediction} evaluated at precision $\beta$, with the current stimulus and category distances held fixed. The feedback likelihood and the category distribution conditional on that feedback are
\begin{align}
\ell_t(h;\beta)&=\sum_{c=1}^{C}a_{t,h}(c)p_{t,h}(c;\beta),\\
p_t^{\mathrm{fb}}(c\mid h)&=
\frac{a_{t,h}(c)p_{t,h}(c)}{\ell_t(h;\beta_t(h))}.
\label{eq:feedback-category-distribution}
\end{align}
The conditional distribution $p_t^{\mathrm{fb}}(c\mid h)$ identifies which category predictions can explain the feedback. It is computed for the current update and is discarded afterwards.

The precision gradient compares the expected category distance before and after conditioning on feedback:
\begin{equation}
\begin{aligned}
\kappa_t(h)
&=\left.\frac{\partial\log\ell_t(h;\beta)}{\partial\beta}\right|_{\beta=\beta_t(h)}\\
&=\sum_{c=1}^{C}p_{t,h}(c)d_{t,h,c}
-\sum_{c=1}^{C}p_t^{\mathrm{fb}}(c\mid h)d_{t,h,c}.
\end{aligned}
\label{eq:precision-gradient}
\end{equation}
Distances and feedback compatibilities remain fixed when taking this derivative. A positive gradient indicates that increasing precision improves the feedback likelihood locally; a negative gradient favors decreasing precision. The gradient is zero when the two expected distances coincide.

A single participant-level learning rate $\eta\in[0,1]$ controls the update in both directions:
\begin{equation}
\beta_{t+1}(h)=\operatorname{clip}
\bigl(\beta_t(h)+\eta\kappa_t(h),\,\beta_{\min},\,\beta_{\max}\bigr).
\label{eq:precision-update}
\end{equation}
Here $\operatorname{clip}(v,l,u)=\min\{u,\max\{l,v\}\}$ restricts precision to the shared bounds $\beta_{\min}=0.1$ and $\beta_{\max}=25$. All active rules update, and the new precisions affect predictions on the next trial. Precision learning has two participant-level parameters: the initial precision $\beta_0$ and the learning rate $\eta$.

Belief and precision update in parallel using the predictions formed before feedback. Belief measures relative support among rules, while precision calibrates the concentration of each rule's category predictions.

\subsection*{Hypothesis updating (H): managing available rules}
\label{m:hypothesis}

Hypothesis updating (H) controls which rules are available for learning. Search responds to recent performance through an accumulated failure state,
\begin{equation}
F_{t+1}=\delta F_t+(1-\delta)(1-r_t),\qquad F_1=0.
\label{eq:failure-accumulator}
\end{equation}
Failure retention $0\leq\delta<1$ sets the persistence of this performance signal. Accumulated failure can increase the probability of searching, $e_{t+1}$, and the probability of a global source conditional on search, $g_{t+1}$:
\begin{align}
e_{t+1}&=\sigmoid[\logit(e_0)+a_eF_{t+1}],\\
g_{t+1}&=g_0+(1-g_0)a_gF_{t+1}.
\label{eq:search-control}
\end{align}
Here $e_0$ and $g_0$ are baseline probabilities, $a_e$ and $a_g$ determine failure sensitivity, $\sigmoid(v)=1/(1+e^{-v})$, and $\logit(p)=\log[p/(1-p)]$. Successful feedback reduces $F$, returning search toward its baseline.

After feedback learning, each rule has exit probability $\rho[1-\pi_t^+(h)]$, where $\rho$ controls the overall exit rate. Weakly supported candidates are therefore more likely to leave, while at least one candidate remains available for prediction and learning.

\label{m:proposals}
Each search event introduces one rule. Local search selects a retained rule according to its belief and generates a related candidate through either a structural edit or functional similarity. Structural edits change one part of a rule, such as the feature used for a judgment. Edit distance measures the cost of this change, and candidates with smaller distances receive greater probability. Functional similarity measures the proportion of reference stimuli assigned the same category by two rules. Candidates with greater agreement receive greater probability. A participant-level mixture weight controls these two local routes. Global search samples uniformly from the catalogue entries outside the current active set. The local and global routes allow hypotheses to change through focused revision and broader exploration.

\label{m:transfer}
If a search event would exceed capacity, the least-supported retained rule is removed. Candidate count therefore varies between one and $M^{\max}$. Retained beliefs are normalized before a share is transferred to the entrant. Transfer depends on functional affinity, which increases with classification agreement. Each old rule divides its belief between itself and the entrant; more strongly supported and more similar rules contribute larger amounts. This allocation preserves a total belief of one. Retained rules preserve their updated precisions, while entrants begin at $\beta_0$. This transfer allows a newly considered rule to build on existing classification knowledge.

\subsection*{Learning hierarchical category relations}
\label{m:pairing}

Graded feedback provides evidence about both classification and family membership. In condition~3, the model maintains joint beliefs $\Pi_t(h,k)$ over rules and the three possible pairings $k$ of four response labels. Pairing beliefs initially assign equal probability to each possibility. Let $\operatorname{mate}_k(y)$ denote the other response in the same family under pairing $k$. The feedback likelihood is
\begin{equation}
\ell_t(h,k)=
\begin{cases}
p_{t,h}(y_t),&r_t=1,\\
p_{t,h}[\operatorname{mate}_k(y_t)],&r_t=0.5,\\
1-p_{t,h}(y_t)-p_{t,h}[\operatorname{mate}_k(y_t)],&r_t=0.
\end{cases}
\label{eq:graded-likelihood}
\end{equation}
Evidence memory applies the discounted update $\Pi_t^+(h,k)\propto[\Pi_t^-(h,k)]^{\gamma}\ell_t(h,k)$, normalized across active rule--pairing combinations. Summing over pairings yields each rule's belief. This joint representation lets partial feedback refine classification as response-pair knowledge develops.

For precision updating, the pairing distribution is evaluated after memory discounting and before incorporating the current feedback:
\begin{equation}
\widetilde q_t^-(k\mid h)=
\frac{[\Pi_t^-(h,k)]^{\gamma}}
{\sum_{k'=1}^{3}[\Pi_t^-(h,k')]^{\gamma}}.
\label{eq:discounted-pairing}
\end{equation}
When $\gamma=0$, active rule--pairing combinations receive equal weight. A rule with zero total joint belief uses a uniform conditional pairing distribution.

For each pairing $k$, define the set of categories compatible with the feedback:
\begin{equation}
\mathcal Y_k(y_t,r_t)=
\begin{cases}
\{y_t\},&r_t=1,\\
\{\operatorname{mate}_k(y_t)\},&r_t=0.5,\\
\{1,2,3,4\}\setminus\{y_t,\operatorname{mate}_k(y_t)\},&r_t=0.
\end{cases}
\label{eq:graded-feedback-categories}
\end{equation}
These sets contain the chosen category, its paired category, or the two categories in the other family, respectively. Averaging over the pairing distribution gives the feedback compatibility
\begin{equation}
a_{t,h}(c)=\sum_{k=1}^{3}\widetilde q_t^-(k\mid h)
\ind\bigl(c\in\mathcal Y_k(y_t,r_t)\bigr).
\label{eq:graded-feedback-compatibility}
\end{equation}
The resulting likelihood equals $\ell_t(h)=\sum_k\widetilde q_t^-(k\mid h)\ell_t(h,k)$. It enters the same precision gradient and update in Eqs.~\eqref{eq:precision-gradient} and~\eqref{eq:precision-update}. Differentiation holds the discounted pairing distribution and the category distances fixed. Thus, all three feedback levels calibrate precision using the pairing knowledge available before that feedback.

During candidate changes, retained rules preserve their conditional pairing beliefs. An entrant inherits the overall pairing distribution obtained after feedback and before candidate removal, carrying learned family relations into the new hypothesis.

\subsection*{Inference from choices}
\label{m:inference}

Particle filtering represents uncertainty over possible internal learning trajectories. Each particle $n$ carries a candidate set and its associated cognitive states. Its normalized weight $W_t^{(n),\mathrm{pre}}$ expresses the researcher's support for that trajectory from preceding observations. The predictive choice probability and weights after observing the current choice are
\begin{align}
\widehat P_t(c)&=\sum_nW_t^{(n),\mathrm{pre}}P_t(c\mid z_t^{(n)}),\\
W_t^{(n),\mathrm{choice}}&=
\frac{W_t^{(n),\mathrm{pre}}P_t(y_t\mid z_t^{(n)})}
{\widehat P_t(y_t)}.
\label{eq:particle-prediction}
\end{align}
Each particle then updates the learner's states using the observed feedback. Within a trajectory, feedback changes the learner's beliefs about rules. Across trajectories, choices change the researcher's support for possible internal states. The likelihood thus averages over uncertainty about perception and candidate dynamics.

The full model defines 13 participant-level parameters, including $\beta_0$ and $\eta$ for precision learning. Free and fixed parameters are specified before estimation. Participant parameters $\bm\theta$ are estimated by minimizing mean choice negative log-likelihood,
\begin{equation}
\mathrm{NLL}(\bm\theta)=-\frac1T\sum_{t\in\mathcal T_{\mathrm{score}}}
\log\widehat P_t(y_t;\bm\theta),
\label{eq:choice-nll}
\end{equation}
where $\mathcal T_{\mathrm{score}}$ is the set of scored trials and $T$ its size. Predictions are formed before current feedback. Choice likelihood is the fitting objective.

\draftslot{Implementation and validation of the revised precision update remain pending. Analysis cohort, estimation settings, training and evaluation boundaries, and uncertainty procedures to be specified.}

\subsection*{Interpreting reconstructed trajectories}
\label{m:measures}

Weighted particle states describe when a rule becomes available, how much belief it receives, how strongly it influences choice and how concentrated its predictions become. These quantities distinguish hypothesis discovery, evidence accumulation and policy expression within individual learning sequences. Candidate-set size measures the number of available explanations; choice-weight concentration measures their effective participation. Their timing can be examined jointly to identify delays between availability, support and dominance. Pre-choice weights summarize states using information available before the current response, and choice-conditioned weights incorporate that response. Analyses specify which conditioning is used and retain the distribution across trajectories as the model's uncertainty about the learner's internal state.

\subsection*{Independent strategy validation}
\label{m:validation}

Verbal reports are reserved for independent evaluation of reconstructed strategies. They are excluded from parameter fitting and particle inference. Validation compares model-predicted strategy summaries with reports encoded according to a prespecified scheme. Comparisons account for the timing of each report and the uncertainty across inferred trajectories.

\draftslot{Report encoding, model-to-report alignment and validation statistics to be specified.}
